\section{Referencial Teórico e Estado da Arte}
\label{sec:referencial}

\subsection{Fundamentos Jurídico-Regulatórios e Triagem Indiciária}

As contratações públicas devem assegurar tratamento isonômico entre os participantes e selecionar a proposta apta a produzir o resultado mais vantajoso para a Administração, conforme os princípios e objetivos estabelecidos pela Lei nº 14.133/2021~\cite{brasil2021licitacoes}. O diploma também prevê infrações administrativas relacionadas à fraude e à frustração do caráter competitivo dos procedimentos licitatórios. No plano concorrencial, a Lei nº 12.529/2011 considera infrações à ordem econômica condutas concertadas entre concorrentes que tenham por objeto ou possam produzir efeitos anticoncorrenciais~\cite{brasil2011concorrencia}. A Lei nº 12.846/2013, por sua vez, estabelece a responsabilização administrativa e civil de pessoas jurídicas pela prática de atos lesivos contra a Administração Pública, incluindo condutas relacionadas a procedimentos licitatórios~\cite{brasil2013anticorrupcao}.

Esse enquadramento jurídico justifica a adoção de mecanismos de prevenção e detecção de sinais de risco, mas não autoriza equiparar uma coincidência cadastral à comprovação de conluio. A existência de sócios em comum, endereços coincidentes, telefones compartilhados ou sobrenomes semelhantes pode constituir elemento para aprofundamento da análise, desde que interpretada em conjunto com o contexto da contratação e com outros elementos objetivos. A abordagem deste trabalho é, portanto, de triagem indiciária: os resultados destinam-se a apoiar a priorização de verificações, não a estabelecer responsabilidade administrativa ou penal.

A proteção de dados pessoais também deve ser considerada. A Lei nº 13.709/2018 (LGPD) disciplina o tratamento de dados pessoais de pessoas naturais e prevê disposições específicas para tratamentos realizados para fins acadêmicos, observadas as condições legais aplicáveis~\cite{brasil2018lgpd}. A pseudonimização reduz a exposição direta de identificadores, mas não equivale necessariamente à anonimização e não afasta, por si só, a incidência da legislação. Além disso, pessoas jurídicas não são titulares de dados pessoais para os fins da LGPD, embora a substituição de seus identificadores por rótulos possa ser adotada como medida de apresentação e governança dos dados. Neste estudo, os identificadores pessoais e empresariais são substituídos por códigos no material divulgado.

\subsection{Redes de Co-disputa, Relações Societárias e Homofilia}

A análise de riscos em licitações pode ser formulada como um problema relacional: além das características individuais dos fornecedores, interessam as relações estabelecidas entre empresas que participam dos mesmos procedimentos e entre essas empresas e seus sócios ou controladores. Guias da Organização para a Cooperação e Desenvolvimento Econômico (OCDE) e do Conselho Administrativo de Defesa Econômica (CADE) descrevem práticas colusivas como divisão de mercado, rodízio de vencedores, propostas de cobertura e supressão de propostas~\cite{oecd2025guidelines,cade2019guia}. Indicadores observáveis nos dados administrativos podem, assim, auxiliar a seleção de procedimentos que mereçam análise aprofundada~\cite{carbone2024bidrigging,fazekas2016objective,wachs2020corruption}.

Estudos sobre redes de co-participação demonstram que as relações entre fornecedores podem revelar padrões que não são evidentes quando cada procedimento é examinado isoladamente~\cite{wachs2019cartel,lyra2021characterization}. Pesquisas sobre propriedade e gestão também destacam a importância de considerar vínculos diretos e indiretos entre empresas concorrentes~\cite{villamil2024collusion,pompeu2025collusion}. Essas abordagens fundamentam a integração, neste trabalho, entre a camada de co-disputa e a camada de propriedade societária.

A interpretação de coincidências nominais exige cautela. O conceito de homofilia descreve a tendência de atores com características semelhantes estabelecerem relações entre si~\cite{mcpherson2001homophily}. Em dados cadastrais, sobrenomes frequentes podem ocorrer em empresas sem relação societária ou operacional. Por esse motivo, a coincidência de sobrenomes é tratada como atributo de menor força do que a identidade direta de sócios, endereço ou telefone. A ponderação discreta adotada neste trabalho é inspirada no princípio de reduzir a influência de atributos frequentes, mas não corresponde ao cálculo clássico de frequência inversa de documentos (IDF).

\subsection{Grafos Heterogêneos e Multicamadas}

Redes heterogêneas permitem representar entidades de diferentes tipos em uma mesma estrutura, enquanto modelos multicamada distinguem categorias de relações que coexistem entre essas entidades~\cite{kivela2014multilayer,newman2018networks}. Essa distinção é relevante para o problema estudado: órgãos públicos, empresas licitantes e sócios ou controladores constituem tipos distintos de vértices; participação em licitação, co-disputa, propriedade societária e vínculo suspeito constituem camadas relacionais distintas.

O grafo utilizado pode ser formalizado por
\begin{equation}
  G=(V,E,\tau_V,\tau_E,w),
  \label{eq:grafo_multicamada}
\end{equation}
em que $V$ é o conjunto de vértices, $E$ é o conjunto de arestas, $\tau_V:V\rightarrow\mathcal{T}_V$ associa cada vértice a um tipo, $\tau_E:E\rightarrow\mathcal{T}_E$ identifica a camada de cada aresta e $w:E\rightarrow\mathbb{R}_{>0}$ atribui pesos às relações.

Neste estudo, os tipos de vértices são órgãos públicos, empresas licitantes e sócios ou controladores. As quatro camadas de arestas são:
\begin{itemize}
  \item \texttt{Participou\_Licitacao}: conecta empresas aos órgãos responsáveis pelos certames dos quais participaram;
  \item \texttt{Concorreu\_Contra}: conecta pares de empresas que disputaram os mesmos certames ou itens, com peso associado à recorrência de co-disputa;
  \item \texttt{É\_Dono\_De}: conecta sócios ou controladores às empresas às quais estão vinculados;
  \item \texttt{Vinculo\_Suspeito}: conecta empresas co-disputantes para as quais foi identificado escore positivo acima do limite mínimo de retenção.
\end{itemize}

As camadas \texttt{Participou\_Licitacao} e \texttt{É\_Dono\_De} são bipartidas entre tipos distintos de vértices. Já \texttt{Concorreu\_Contra} e \texttt{Vinculo\_Suspeito} conectam empresas entre si. Consequentemente, o grafo agregado não é, em geral, bipartido, ainda que algumas de suas camadas o sejam. Essa distinção é necessária para interpretar corretamente os motivos triangulares examinados neste trabalho.

\subsection{Motivos Triangulares e Compartilhamento Societário}

Um triângulo $K_3$ é um subgrafo com três vértices distintos e uma aresta entre cada par de vértices. No problema estudado, considera-se um motivo triangular tipado formado por duas relações societárias e uma relação entre empresas concorrentes. Para um sócio $s$ compartilhado pelas empresas $a$ e $b$, o motivo é
\begin{equation}
  \{(s,a),(s,b),(a,b)\},
  \label{eq:motivo_triangular}
\end{equation}
em que $(s,a)$ e $(s,b)$ pertencem à camada \texttt{É\_Dono\_De}, enquanto $(a,b)$ pertence à camada \texttt{Vinculo\_Suspeito}.

O triângulo, portanto, não está presente na camada bipartida de propriedade considerada isoladamente. Ele aparece na representação agregada que combina as duas relações de propriedade com a aresta entre as empresas. Se duas empresas concorrentes compartilham dois sócios ou controladores distintos, cada elemento da interseção dos conjuntos de sócios forma um motivo triangular. A existência desses motivos caracteriza uma sobreposição societária observável, mas não comprova, por si só, coordenação ilícita ou ausência efetiva de independência decisória.

\subsection{Fundamentos dos Algoritmos de Redução Topológica}

A rede completa contém relações de natureza distinta e vértices com diferentes níveis de conectividade. Para concentrar a análise nas relações prioritárias, aplica-se um predicado de retenção às arestas e, em seguida, removem-se os vértices que não possuem incidência no subgrafo filtrado. Essa operação deve ser distinguida de um subgrafo induzido por conjunto de vértices: o procedimento implementado seleciona arestas por tipo e peso e preserva inicialmente o conjunto de vértices, inclusive os que se tornam isolados.

O $1$-core de um grafo é o maior subgrafo induzido em que todos os vértices possuem grau mínimo igual ou superior a um. Para $k=1$, sua obtenção equivale à remoção dos vértices isolados, isto é, aqueles cujo grau é zero. Como esses vértices não possuem arestas incidentes, sua remoção não provoca uma cascata de exclusões em outros vértices. A noção de $k$-core é utilizada na análise de coesão estrutural de redes~\cite{seidman1983network}. Quando a operação consiste apenas em remover vértices de grau zero, a descrição mais precisa é remoção de isolados para obtenção do $1$-core.

Após a filtragem, a identificação de componentes conexas separa os conjuntos de vértices que permanecem conectados por caminhos. A estrutura disjoint-set, também denominada Union-Find, mantém uma partição dinâmica dos vértices por meio de operações de busca e união. Com compressão de caminho e união por posto ou tamanho, o custo amortizado de uma sequência de operações é quase linear, podendo ser expresso por
\begin{equation}
  O\big((|V|+|E|)\alpha(|V|)\big),
  \label{eq:complexidade_union_find}
\end{equation}
em que $\alpha$ é a função inversa de Ackermann~\cite{tarjan1975efficiency}.

A enumeração dos motivos triangulares pode explorar a estrutura específica do problema. Para cada aresta suspeita $(a,b)$ entre empresas, comparam-se os conjuntos de sócios associados a cada empresa:
\begin{equation}
  \mathrm{Donos}(a)\cap\mathrm{Donos}(b).
  \label{eq:intersecao_donos}
\end{equation}
Cada sócio presente nessa interseção forma, com as duas empresas, um motivo triangular tipado. Com conjuntos previamente indexados, a operação pode ser implementada por interseção de conjuntos; seu custo depende dos tamanhos dos conjuntos comparados. Uma estimativa superior simplificada é $O(|E_s|d_{\max})$, em que $E_s$ é o conjunto de arestas suspeitas examinadas e $d_{\max}$ é o maior número de sócios associado a uma empresa, assumindo consultas de pertinência em tempo médio constante.

\subsection{Ponderação de Similaridade e Controle de Falsos Positivos}

O coeficiente de Jaccard é uma medida de similaridade entre conjuntos, definida por
\begin{equation}
  J(A,B)=\frac{|A\cap B|}{|A\cup B|}.
\end{equation}
Embora útil em diversas aplicações, a normalização pela união pode reduzir a expressão da quantidade absoluta de itens compartilhados quando os conjuntos comparados possuem tamanhos muito diferentes. Por essa razão, a metodologia preserva contagens absolutas de co-disputa como atributo de recorrência, sem incorporá-las diretamente ao escore cadastral e societário.

A frequência inversa de documentos (IDF) reduz a contribuição de termos frequentes em uma coleção. Uma formulação clássica é
\begin{equation}
  \mathrm{IDF}(t)=\log\left(\frac{N}{f(t)}\right),
\end{equation}
em que $N$ representa o número de documentos e $f(t)$ a quantidade de documentos que contêm o termo $t$. A implementação deste estudo não calcula essa expressão. Em seu lugar, utiliza uma ponderação discreta: sobrenomes previamente classificados como triviais recebem peso $0{,}25$, enquanto os demais recebem peso $1{,}00$. A regra é inspirada no princípio de reduzir a influência de atributos frequentes, mas deve ser caracterizada como heurística determinística, e não como aplicação integral de IDF.

\subsection{Síntese do Enquadramento e Limites Interpretativos}

O referencial sustenta a integração de relações licitatórias, societárias e cadastrais em um grafo heterogêneo multicamada. A literatura de redes fornece instrumentos para representar e examinar essas relações; os fundamentos algorítmicos justificam a filtragem e a redução estrutural; e o enquadramento jurídico delimita os resultados como sinais para triagem e priorização de auditorias. A metodologia não busca inferir automaticamente a ocorrência de conluio, mas identificar configurações relacionais que mereçam verificação documental e contextual pelos órgãos competentes.